\documentclass[a4paper]{article}

\def\npart{III}

\def\ntitle{Mapping Class Groups}
\def\nlecturer{H.\ Wilton}

\def\nterm{Michaelmas}
\def\nyear{2019}

\ifx \nauthor\undefined
  \def\nauthor{Qiangru Kuang}
\else
\fi

\ifx \ntitle\undefined
  \def\ntitle{Template}
\else
\fi

\ifx \nauthoremail\undefined
  \def\nauthoremail{qk206@cam.ac.uk}
\else
\fi

\ifx \ndate\undefined
  \def\ndate{\today}
\else
\fi

\title{\ntitle}
\author{\nauthor}
\date{\ndate}

%\usepackage{microtype}
\usepackage{mathtools}
\usepackage{amsthm}
\usepackage{stmaryrd}%symbols used so far: \mapsfrom
\usepackage{empheq}
\usepackage{amssymb}
\let\mathbbalt\mathbb
\let\pitchforkold\pitchfork
\usepackage{unicode-math}
\let\mathbb\mathbbalt%reset to original \mathbb
\let\pitchfork\pitchforkold

\usepackage{imakeidx}
\makeindex[intoc]

%to address the problem that Latin modern doesn't have unicode support for setminus
%https://tex.stackexchange.com/a/55205/26707
\AtBeginDocument{\renewcommand*{\setminus}{\mathbin{\backslash}}}
\AtBeginDocument{\renewcommand*{\models}{\vDash}}%for \vDash is same size as \vdash but orginal \models is larger
\AtBeginDocument{\let\Re\relax}
\AtBeginDocument{\let\Im\relax}
\AtBeginDocument{\DeclareMathOperator{\Re}{Re}}
\AtBeginDocument{\DeclareMathOperator{\Im}{Im}}
\AtBeginDocument{\let\div\relax}
\AtBeginDocument{\DeclareMathOperator{\div}{div}}

\usepackage{tikz}
\usetikzlibrary{automata,positioning}
\usepackage{pgfplots}
%some preset styles
\pgfplotsset{compat=1.15}
\pgfplotsset{centre/.append style={axis x line=middle, axis y line=middle, xlabel={$x$}, ylabel={$y$}, axis equal}}
\usepackage{tikz-cd}
\usepackage{graphicx}
\usepackage{newunicodechar}

\usepackage{fancyhdr}

\fancypagestyle{mypagestyle}{
    \fancyhf{}
    \lhead{\emph{\nouppercase{\leftmark}}}
    \rhead{}
    \cfoot{\thepage}
}
\pagestyle{mypagestyle}

\usepackage{titlesec}
\newcommand{\sectionbreak}{\clearpage} % clear page after each section
\usepackage[perpage]{footmisc}
\usepackage{blindtext}

%\reallywidehat
%https://tex.stackexchange.com/a/101136/26707
\usepackage{scalerel,stackengine}
\stackMath
\newcommand\reallywidehat[1]{%
\savestack{\tmpbox}{\stretchto{%
  \scaleto{%
    \scalerel*[\widthof{\ensuremath{#1}}]{\kern-.6pt\bigwedge\kern-.6pt}%
    {\rule[-\textheight/2]{1ex}{\textheight}}%WIDTH-LIMITED BIG WEDGE
  }{\textheight}% 
}{0.5ex}}%
\stackon[1pt]{#1}{\tmpbox}%
}

%\usepackage{braket}
\usepackage{thmtools}%restate theorem
\usepackage{hyperref}

% https://en.wikibooks.org/wiki/LaTeX/Hyperlinks
\hypersetup{
    %bookmarks=true,
    unicode=true,
    pdftitle={\ntitle},
    pdfauthor={\nauthor},
    pdfsubject={Mathematics},
    pdfcreator={\nauthor},
    pdfproducer={\nauthor},
    pdfkeywords={math maths \ntitle},
    colorlinks=true,
    linkcolor={red!50!black},
    citecolor={blue!50!black},
    urlcolor={blue!80!black}
}

\usepackage{cleveref}



% TODO: mdframed often gives bad breaks that cause empty lines. Would like to switch to tcolorbox.
% The current workaround is to set innerbottommargin=0pt.

%\usepackage[theorems]{tcolorbox}





\usepackage[framemethod=tikz]{mdframed}
\mdfdefinestyle{leftbar}{
  %nobreak=true, %dirty hack
  linewidth=1.5pt,
  linecolor=gray,
  hidealllines=true,
  leftline=true,
  leftmargin=0pt,
  innerleftmargin=5pt,
  innerrightmargin=10pt,
  innertopmargin=-5pt,
  % innerbottommargin=5pt, % original
  innerbottommargin=0pt, % temporary hack 
}
%\newmdtheoremenv[style=leftbar]{theorem}{Theorem}[section]
%\newmdtheoremenv[style=leftbar]{proposition}[theorem]{proposition}
%\newmdtheoremenv[style=leftbar]{lemma}[theorem]{Lemma}
%\newmdtheoremenv[style=leftbar]{corollary}[theorem]{corollary}

\newtheorem{theorem}{Theorem}[section]
\newtheorem{proposition}[theorem]{Proposition}
\newtheorem{lemma}[theorem]{Lemma}
\newtheorem{corollary}[theorem]{Corollary}
\newtheorem{axiom}[theorem]{Axiom}
\newtheorem*{axiom*}{Axiom}

\surroundwithmdframed[style=leftbar]{theorem}
\surroundwithmdframed[style=leftbar]{proposition}
\surroundwithmdframed[style=leftbar]{lemma}
\surroundwithmdframed[style=leftbar]{corollary}
\surroundwithmdframed[style=leftbar]{axiom}
\surroundwithmdframed[style=leftbar]{axiom*}

\theoremstyle{definition}

\newtheorem*{definition}{Definition}
\surroundwithmdframed[style=leftbar]{definition}

\newtheorem*{slogan}{Slogan}
\newtheorem*{eg}{Example}
\newtheorem*{ex}{Exercise}
\newtheorem*{remark}{Remark}
\newtheorem*{notation}{Notation}
\newtheorem*{convention}{Convention}
\newtheorem*{assumption}{Assumption}
\newtheorem*{question}{Question}
\newtheorem*{answer}{Answer}
\newtheorem*{note}{Note}
\newtheorem*{application}{Application}

%operator macros

%basic
\DeclareMathOperator{\lcm}{lcm}

%matrix
\DeclareMathOperator{\tr}{tr}
\DeclareMathOperator{\Tr}{Tr}
\DeclareMathOperator{\adj}{adj}

%algebra
\DeclareMathOperator{\Hom}{Hom}
\DeclareMathOperator{\End}{End}
\DeclareMathOperator{\id}{id}
\DeclareMathOperator{\im}{im}
\DeclarePairedDelimiter{\generation}{\langle}{\rangle}

%groups
\DeclareMathOperator{\sym}{Sym}
\DeclareMathOperator{\sgn}{sgn}
\DeclareMathOperator{\inn}{Inn}
\DeclareMathOperator{\aut}{Aut}
\DeclareMathOperator{\GL}{GL}
\DeclareMathOperator{\SL}{SL}
\DeclareMathOperator{\PGL}{PGL}
\DeclareMathOperator{\PSL}{PSL}
\DeclareMathOperator{\SU}{SU}
\DeclareMathOperator{\UU}{U}
\DeclareMathOperator{\SO}{SO}
\DeclareMathOperator{\OO}{O}
\DeclareMathOperator{\PSU}{PSU}

%hyperbolic
\DeclareMathOperator{\sech}{sech}

%field, galois heory
\DeclareMathOperator{\ch}{ch}
\DeclareMathOperator{\gal}{Gal}
\DeclareMathOperator{\emb}{Emb}



%ceiling and floor
%https://tex.stackexchange.com/a/118217/26707
\DeclarePairedDelimiter\ceil{\lceil}{\rceil}
\DeclarePairedDelimiter\floor{\lfloor}{\rfloor}


\DeclarePairedDelimiter{\innerproduct}{\langle}{\rangle}

%\DeclarePairedDelimiterX{\norm}[1]{\lVert}{\rVert}{#1}
\DeclarePairedDelimiter{\norm}{\lVert}{\rVert}



%Dirac notation
%TODO: rewrite for variable number of arguments
\DeclarePairedDelimiterX{\braket}[2]{\langle}{\rangle}{#1 \delimsize\vert #2}
\DeclarePairedDelimiterX{\braketthree}[3]{\langle}{\rangle}{#1 \delimsize\vert #2 \delimsize\vert #3}

\DeclarePairedDelimiter{\bra}{\langle}{\rvert}
\DeclarePairedDelimiter{\ket}{\lvert}{\rangle}




%macros

%general

%divide, not divide
\newcommand*{\divides}{\mid}
\newcommand*{\ndivides}{\nmid}
%vector, i.e. mathbf
%https://tex.stackexchange.com/a/45746/26707
\newcommand*{\V}[1]{{\ensuremath{\symbf{#1}}}}
%closure
\newcommand*{\cl}[1]{\overline{#1}}
%conjugate
\newcommand*{\conj}[1]{\overline{#1}}
%set complement
\newcommand*{\stcomp}[1]{\overline{#1}}
\newcommand*{\compose}{\circ}
\newcommand*{\nto}{\nrightarrow}
\newcommand*{\p}{\partial}
%embed
\newcommand*{\embed}{\hookrightarrow}
%surjection
\newcommand*{\surj}{\twoheadrightarrow}
%power set
\newcommand*{\powerset}{\mathcal{P}}

%matrix
\newcommand*{\matrixring}{\mathcal{M}}

%groups
\newcommand*{\normal}{\trianglelefteq}
%rings
\newcommand*{\ideal}{\trianglelefteq}

%fields
\renewcommand*{\C}{{\mathbb{C}}}
\newcommand*{\R}{{\mathbb{R}}}
\newcommand*{\Q}{{\mathbb{Q}}}
\newcommand*{\Z}{{\mathbb{Z}}}
\newcommand*{\N}{{\mathbb{N}}}
\newcommand*{\F}{{\mathbb{F}}}
%not really but I think this belongs here
\newcommand*{\A}{{\mathbb{A}}}

%asymptotic
\newcommand*{\bigO}{O}
\newcommand*{\smallo}{o}

%probability
\newcommand*{\prob}{\mathbb{P}}
\newcommand*{\E}{\mathbb{E}}

%vector calculus
\newcommand*{\gradient}{\V \nabla}
\newcommand*{\divergence}{\gradient \cdot}
\newcommand*{\curl}{\gradient \cdot}

%logic
\newcommand*{\yields}{\vdash}
\newcommand*{\nyields}{\nvdash}

%differential geometry
\renewcommand*{\H}{\mathbb{H}}
\newcommand*{\transversal}{\pitchfork}
\renewcommand{\d}{\mathrm{d}} % exterior derivative

%number theory
\newcommand*{\legendre}[2]{\genfrac{(}{)}{}{}{#1}{#2}}%Legendre symbol


\graphicspath{{figures/}}

\DeclareMathOperator{\Homeo}{Homeo}
\DeclareMathOperator{\Diff}{Diff}
\DeclareMathOperator{\Mod}{Mod}
\renewcommand{\b}{\p}
\DeclareMathOperator{\Isom}{Isom}
\DeclareMathOperator{\Fix}{Fix}
\DeclareMathOperator{\Axis}{Axis}

\begin{document}

\begin{titlepage}
  \begin{center}
    \includegraphics[width=0.6\textwidth]{logo.jpg}\par
    \vspace{1cm}
    {\scshape\huge Mathematics Tripos \par}
    \vspace{2cm}
    {\huge Part \npart \par}
    \vspace{0.6cm}
    {\Huge \bfseries \ntitle \par}
    \vspace{1.2cm}
    {\Large\nterm, \nyear \par}
    \vspace{2cm}
    
    {\large \emph{Lectures by } \par}
    \vspace{0.2cm}
    {\Large \scshape \nlecturer}
    
    \vspace{0.5cm}
    {\large \emph{Notes by }\par}
    \vspace{0.2cm}
    {\Large \scshape \href{mailto:\nauthoremail}{\nauthor}}
 \end{center}
\end{titlepage}

\tableofcontents

\setcounter{section}{-1}

\section{Introduction}

\subsection{Classification of surfaces}

The subject of this course is surfaces, i.e.\ two-dimensional manifolds. We assume throughout the manifold \(S\) is connected, smooth and oriented. More importantly, we assume it is of \emph{finite type}: \(S = \overline S - \{\text{finite set}\}\) where \(\overline S\) is a compact manifold possibly with boundary, and the finite set is contained in its interior.

\begin{theorem}[classification of surfaces of finite type]
  Every connected orientable surface of finite type is diffeomorphic to some \(S_{g, n, b}\) which is a surface of genus \(g\) with \(n\) punctures and \(b\) boundary components.
\end{theorem}

\begin{figure}[ht]
  \centering
  \def\svgwidth{.8\columnwidth}
  \input{figures/surface.pdf_tex}
  \caption{\(S_{g, n, b}\)}
\end{figure}

It has Euler characteristic \(\chi(S) = 2 - 2g - (n + b)\).

A \emph{closed surface} is a surface with \(n = b = 0\). A \emph{compact surface} is a surface with \(n = 0\).

\begin{eg}
  Suppose \(\chi(S) > 0\) then \(g = 0, n + 1 = 0 \text{ or } 1\). Thus \(S\) is either \(S^2\), or \(S^2\) with a puncture which is isomorphic to \(\C\), or \(S^2\) with one boundary component which is isomorphic of \(D^2\), which can be thought as sitting inside \(\C\) as the unit disk.

  Note that we can either think of a puncture as a deleted point, or as a marked point on the surface. These two views are equivalent and we will use whichever that is more convenient.
\end{eg}

\begin{eg}
  Suppose \(\chi(S) = 0\) then either \(g = 1, n + b = 0\) or \(g = 0, n + b = 2\). Thus it is either the torus, the punctured plane \(\C_*\), \(S^1 \times I\) or the punctured disk \(D^2_*\).
\end{eg}

\subsection{Mapping class groups}

The natural group associated to a surface \(S\) is the group of orientation-preserving homeomorphism \(\Homeo^+(S)\). However this is a huge group. It might be easier to consider the maps up to homotopy. Suppose \(\phi_0, \phi_1: S \to S\) are homeomorphisms. We can define an equivalence relation \(\phi_0 \sim \phi_1\) if there exists an isotopy \(\phi_t: S \times I \to S\) from \(\phi_0\) to \(\phi_1\).

Another way to think about this: give \(\Homeo^+(S)\) the compact-open topology. Let \(\Homeo_0(S)\) be the path component of \(1_S\). It is an exercise to show that \(\Homeo_0(S)\) is a normal subgroup of \(\Homeo(S)\).

\begin{definition}[mapping class group]\index{mapping class group}
  The \emph{mapping class group} of \(S\) is
  \[
    \Mod(S) = \Homeo^+(S, \b S)/ \Homeo_0(S, \b S)
  \]
  where \(\Homeo(S, \b S)\) is the subgroup of \(\Homeo(S)\) that fixes \(\b S\) pointwise.
\end{definition}

We might ask if we can replace homeomorphism by diffeomorphism, or we replace isotopy by homotopy. Indeed we have

\begin{theorem}[Baer, Munkres]
  For any smooth surface \(S\) of finite type,
  \[
    \Mod(S) \cong \Diff^+(S, \b S)/ \Diff_0(S, \b S) \cong \Homeo^+(S, \b S)/ \sim
  \]
  where the equivalence relation is homotopy.
\end{theorem}

\subsection{Context \& Motivation}

\paragraph{bundles}

Let \(\phi \in \Diff(S)\). We can define \(M_\phi = S \times [0, 1]/ \sim\) where \((x, 1) \sim (\phi(x), 1)\). This is called a \emph{surface bundle} over \(S^1\). Note that \(M_\phi\) only depends on \([\phi] \in \Mod(S)\).

More generally, if \(B\) is a space and \(\rho: \pi_1(B) \to \Mod (S)\) then we get a bundle \(\tilde B \times S\) ... this leads to an \(S\)-bundle over \(B\).

\paragraph{moduli space}

More handwavy motivation. Let \(S = S_g = S_{g, 0, 0}\). It turns out the moduli space \(\mathcal M_g\) of geometric structures on \(S_g\) is the same as the moduli space of complex structures on \(S_g\). Morally (but not actually true), the universal cover \(\tilde{\mathcal M_g}\) is \(\mathcal T_g\), the Teichüller space, and \(\pi_1(\mathcal M_g)\) is \(\Mod(S)\), so we have \(\mathcal M_g = \Mod(S_g) \setminus \mathcal T_g\).

\paragraph{analogy}

There is an analogy between the torus and a surface \(S\).

\begin{table}[h!]
  \centering
  \begin{tabular}{c|c}
    \(S\) & \(T\) \\ \hline
    \(\pi_1(S)\) & \(\Z^n\) \\ \hline
    \(\Mod(S)\) & \(\SL_n(\Z)\) \\ \hline
    closed curves (up to isotopy) & vectors
  \end{tabular}
  \caption{Comparison of \(S\) with \(T^n\)}
\end{table}

\section{Curves, Surfaces \& Hyperbolic geometry}

\subsection{The hyperblic plane}

There are two models of hyperbolic geometry.

\paragraph{The upper half-plane model}

It is the half-plane \(\H^2 = \{x + iy: y > 0\}\) with metric
\[
  ds^2 = \frac{dx^2 + dy^2}{y^2}.
\]
The geodesics are either vertical lines or semicircles on the \(x\)-axis. Then both meet the \(x\)-axis at right angle. The orientation preserving isometries are
\[
  \Isom^+(\H^2) = PSL_2(\R)
\]
where an element in \(PSL_2(\R)\) acts by Möbius transformation:
\[
  \begin{pmatrix}
    a & b \\
    c & d
  \end{pmatrix}
  z = \frac{az + b}{cz + d}
\]
where \(a, b, c, d \in \R\) and \(ad - bc > 0\) (by rescaling). As the entries are real, it preserves the \(x\)-axis and the point at infinity.

\paragraph{The Poincaré disc model}

Conjugating the isometry group by \(z \mapsto \frac{z - i}{z + i}\), we can map \(\H^2\) to the interior unit disc. The metric is
\[
  ds^2 = 4 \frac{dx^2 + dy^2}{(1 - r^2)^2}.
\]
Note that it is radically symmetric. As Möbius transformations are conformal, the geodesics are diameters and arcs meeting the boundary of the disc at right angle.

The (Gromov) boundary (at infinity) \(\b \H^2\) is the unit circle bounding the disc. We write \(\overline \H^2 = \H^2 \cup \b \H^2\).

\begin{note}
  Each isometry \(f\) of \(\H^2\) extend uniquely to Möbius transformation \(\overline f\) of \(\b \H^2\).
\end{note}

Let \(f \in \Isom^+ \H^2 = \PSL_2(\R)\) and let \(n\) be the number of fixed points of \(\overline f\). Brouwer fixed point theorem says \(n \geq 1\). On the other hand, \(n \leq 2\) unless \(f = \id\).

\paragraph{\(n = 2\)}

Let be \(\{\xi_+, \xi_-\} = \Fix(\overline f) \subseteq \overline \H^2\). If one of them is in the interior then there is a geodesic between them lying in \(\H^2\). But then \(f\) fixes every point on the geodesic, absurd. Thus \(\xi_+, \xi_- \in \b \H^2\) and there is a unique geodesic between them.

There exists \(g \in \PSL_2(\R)\) such that \(g(\xi_-) = 0, g(\xi_+) = \infty\) in the upper half-plane model, and \(\Fix(g \overline f g^{-1}) = g \Fix(\overline f)\). Thus after conjugation, wlog \(\xi_- = 0, \xi_+ = \infty\) so
\[
  f(z) = \lambda z = e^{\tau(f)} z
\]
where \(\lambda > 0\). So \(f\) now acts as translation on the \(y\)-axis by \(\tau(f)\).

In general, \(f\) preserves a geodesic line called \(\Axis(f)\), acting by translation. Such an \(f\) is called \emph{hyperbolic} or \emph{loxodromic}.

fact: If \(x \notin \Axis(f)\) then \(d(x, f(x)) > \tau(f)\).

\paragraph{\(n = 1\)}

Let \(\Fix(\overline f) = \{\xi\}\). If \(\xi \in \H^2\) wlog \(\xi = 0\) in the disc model. Then must have \(f(z) = e^{i\theta} z\). Such an \(f\) is called \emph{elliptic}. If \(\xi \in \b \H^2\) wlog \(\xi = \infty\) in the upper half-plane model. It follows that \(f(z) = z \pm 1\). Such an \(f\) is called \emph{parabolic}.

\begin{remark}
  This classification is invariant under conjugacy.
\end{remark}

\subsection{Hyperbolic structures}

A \emph{hyperbolic structure} on \(S\) is a complete, finite-area Riemannian metric of constant curvature \(\kappa = +1, 0 \text{ or } -1\), in which every boundary components are geodesic.

What kind of hyperbolic structure can we put on a surface? Recall Gauss-Bonnet which says that if \(S\) has finite area then
\[
  \int_S \kappa dA = 2\pi \chi(S)
\]
so the sign of \(\kappa\) is the same as the sign of \(\chi(S)\). In particular
\begin{enumerate}
\item if \(\chi(S) > 0\) then \(\kappa = 1\) so \(S\) is locally \(S^2\).
\item if \(\chi(S) = 0\) then \(\kappa = 0\) so \(S\) is locally \(\R^2\).
\item if \(\chi(S) < 0\) then \(\kappa = -1\) so \(S\) is locally \(\H^2\).
\end{enumerate}

\begin{eg}
  Recall that \(\chi(S) > 0\) implies \(S\) is \(S^2, D^2\) or \(\C\). \(S^2\) will just have the geometric structure of the sphere. For \(\C\), there is no complete finite-area metric. For \(D^2\), recall that we require the boundary component to be a geodesic so the geometric structure on \(D^2\) will just be the semisphere.

  If \(\chi(S) = 0\) then \(S\) is \(T^2, A = S^1 \times [0, 1], \C_*\) or \(D_*^2\). Again \(\C_*\) and \(D_*^2\) do not admit finite-area complete metric. \(T\) will then be the flat torus, i.e.\ \(\Z^2 \setminus \R^2\). \(A\) has the geometric structure of a cylinder.
\end{eg}

\begin{theorem}
  If \(S\) is connected, oriented and of finite type and \(\chi(S) < 0\) then there is a convex subspace \(\tilde S \subseteq \H^2\) with geodesic boundary and an action \(\pi_1S\) on \(\tilde S\) by isometry such that
  \[
    S \cong \pi_1 S \setminus \tilde S
  \]
  has finite area. In particular, \(S\) has curvature \(-1\) everywhere.
\end{theorem}

\begin{proof}[Sketch proof when \(S = S_{g, 0, 0}\)]
  Cut along \(2g\) loops to get a \(4g\)-gon with sides identified appropriately. Then it suffices to find a \(4g\)-gon in the hyperbolic disk whose internal angles sum up to \(2\pi\). A regular \(4g\)-gon with vertices on \(\b \H\) has total interior angle \(0\), while as we shrink the polygon it resembles more and more like a Euclidean polygon so the total interior angle approachs \((4g - 2)\pi\). As \(g > 1\), by intermediate value theorem we can find a regular \(4g\)-gon whose total interior angle is \(2\pi\).

  The metric on \(S\) then has constant curvature \(\kappa = -1\). Then by Jacobi theorem \(\tilde S \cong \H^2\) (homeo?). The statement about fundamental group follows from algebraic topology (note that \(\pi_1\) does act by isometry since the upstairs metric is lifted from the downstairs').
\end{proof}

Such a surface \(S\) is called \emph{hyperbolic}.

\begin{remark}
  When \(S\) has no boundary components \(\tilde S = \H^2\).
\end{remark}

\subsection{Curves on hyperbolic surface}

A \emph{closed curve} on \(S\) is a smooth map \(\alpha: S^1 \to S\). It gives a conjugacy class \([\sigma] \in \pi_1S\). This leads to an isometry of \(\H^2\) (up to isometry) when \(S\) is hyperbolic. Thus we could talk about conjugacy-invariant properties of \(\alpha\).

\begin{definition}
  We say \(\alpha\) is \emph{inessential} if \(\alpha\) is homotopic to a point or a puncture. Otherwise \(\alpha\) is essential.
\end{definition}

Picture of an (embedded) hyperbolic surface. Note that a puncture is a cusp because of finite-area and completeness.

\begin{lemma}\leavevmode
  \begin{enumerate}
  \item If \(\alpha\) is elliptic then it is homotopic to a point.
  \item If \(\alpha\) is parabolic then it is homotopic to a puncture.
  \item If \(\alpha\) is hyperbolic then it is essential.
  \end{enumerate}
\end{lemma}

\begin{proof}\leavevmode
  \begin{enumerate}
  \item \(\alpha\) is elliptic implies that \(\alpha\) fixes a point of \(\H^2\). But the action of \(\pi_1\) is free so \(\alpha = \id\) as isometry. so \(\alpha\) is homotopic to a point.
  \item If \(\alpha\) is parabolic then wlog \(\alpha: z \mapsto z + 1\). Take \(\alpha(0) = x_0 \in S\) to be a basepoint and choose \(\tilde x_0\) a lift in \(\H^2\). Let \(\tilde \alpha\) be the lift of \(\alpha\) at \(\tilde x_0\). Let \(\tilde \alpha_s(t) = \tilde \alpha(t) + is\) for \(s \in [0, \infty)\). For all \(s\),
    \[
      \tilde \alpha_s(1) = \tilde \alpha_s(0) + 1
    \]
    so \(\tilde \alpha_s\) descends to a loop \(\alpha_s\) in \(S\) and \(\alpha_s\) tends to a puncture of \(s\) by compactness of \(\overline \H^2\).
  \item It's enough to prove that \(\alpha\) homotopic to a puncture then it is parabolic. By shrinking \(\alpha\) we get a sequence of annuli. Since the hyperbolic structure is complete and of finite area, there exists \(\alpha_n \sim \alpha\) such that \(\ell(\alpha_n) \to 0\) as \(n \to \infty\). Lift \(\alpha\) to a path \(\tilde \alpha: [0, 1] \to \H^2\) and moreover we get well-defined lifts \(\tilde \alpha_n\) of \(\alpha_n\). Let \(\tilde x_n = \tilde \alpha_n(0)\). Note \(\tilde \alpha_n(1) = \alpha \cdot \tilde x_n\). Then the translation distance is
    \begin{align*}
      \tau(\alpha)
      &\leq d(\tilde x_n, \alpha \cdot \tilde x_n) \\
      &= d(\tilde \alpha_n(0), \tilde \alpha_n(1)) \\
      &\leq \ell(\tilde \alpha_n) \\
      &= \ell(\alpha_n) \\
      &\to 0
    \end{align*}
    as \(n \to \infty\) so \(\alpha\) is not hyperbolic.
  \end{enumerate}
\end{proof}

\begin{lemma}
  Let \(S\) be a hyperbolic surface and \(\alpha\) an essential closed curve on \(S\). Then there is a unique geodesic representative in the homotopy class of \(\alpha\).
\end{lemma}

Note that for Euclidean space, such as a torus, such a representative exists but is not unique.

\begin{proof}
  The universal cover of \(S^1\) is \(\R\) and the universal cover of \(\widetilde S \subseteq \H^2\) of \(S\). Then each \(\alpha: S^1 \to S\) lifts to \(\tilde \alpha: \R \to \widetilde S\). Note that \(\tilde \alpha\) is \(\Z\)-equivariant by considering the action of \(\pi_1 S^1\). From the lemma above we know \(\alpha\) has an axis \(\Axis(\alpha)\). Let \(\pi: \H^2 \to \Axis(\alpha)\) be the orthogonal projection (equivalently, it projects a point \(x\) to \(\pi(x)\) on \(\Axis(\alpha)\) such that the length of the geodesic between \(x\) and \(\pi(x)\) is the shortest so. See example sheet 1). Let \(\tilde \gamma_t: [0, 1] \to \H^2\) be the unique constant speed geodesic from \(\tilde \alpha(t)\) to \(\pi \compose \tilde \alpha(t)\). Since \(\langle \alpha \rangle\) acts on both \(\tilde \alpha\) and \(\Axis(\alpha)\) and the \(\tilde \gamma_t\) are determined canonically, taking the quotient by \(\Z = \langle \alpha \rangle\) defines a homotopy from \(\alpha\) to some closed curve \(\beta\) on \(S\) in the image of \(\Axis(\alpha)\). The image of \(\beta\) is a local geodesic and after reparameterisation, \(\beta\) is a constant-speed geodesic.

  Unqiueness: suppose \(\alpha \simeq \beta\) are both geodesics on \(S\). Lift \(\alpha, \beta\) to \(\tilde \alpha, \tilde: \R \to \H^2\) geodesics in \(\H^2\) that are contained in a bounded neighbourhood of each other (since \(\sup_{t \in S^1} d(\alpha(t), \beta(t)) < \infty\) by compactness). It follows that \(\tilde \alpha, \tilde \beta\) are geodesics in \(\H^2\) with the same endpoints on \(\b \H^2\). It follws that \(\tilde \alpha = \tilde \beta\) so \(\alpha = \beta\).
\end{proof}

\section{Simple closed curves \& Intersection number}

A closed curve \(\alpha: S^1 \to S\) is \emph{simple} if it is injective. The idea is that simple closed curves are like basis of a vector space and we can understand mapping class groups by understanding its action on simplex closed curves.

\begin{definition}[(ambient) isotopy of simple closed curves]\index{ambient isotopy}
  A homotopy \(\alpha_\bullet\) between simple closed curves \(\alpha_0\) to \(\alpha_1\) is an \emph{isotopy} if each \(\alpha_t\) is simple.

  If \(\phi_|bullet: S \to S\) is an isotopy such that \(\phi_0 = \id_S\) and \(\phi_1 \compose \alpha_0 = \alpha_1\) then we say \(\alpha_0, \alpha_1\) are \emph{ambient isotopic}.
\end{definition}

\begin{lemma}
  Two essential simple closed curves on an orientable surface \(S\) are homotopic relative to \(\b S\) if and only if they are ambient isotopic.
\end{lemma}

We'll prove this later.

\begin{definition}
  An element \(h \in \pi_1(S)\) is \emph{primitive} if \(h \neq g^n\) for some \(n > 1\).
\end{definition}

\begin{lemma}
  Let \(T^2\) be the torus. The homotopy class of essential simple closed curves on \(T^2\) correspond to primitive elements of \(\pi_1 T^2 = \Z^2\).
\end{lemma}

\begin{proof}
  Example sheet 1, question 8.
\end{proof}

\begin{lemma}
  If \(\alpha\) is an essential simple closed curve on a hyperbolic surface \(S\) then \(\alpha \in \pi_1S\) is primitive. In fact the it has centraliser \(C(\alpha) = \langle \alpha\rangle\).
\end{lemma}

\begin{proof}
  wlog \(\alpha\) is a geodesic and we may consider \(\Axis(\alpha) \subseteq \H^2\). Let \(g \in C(\alpha), x \in \Axis(\alpha)\). Then
  \[
    d(gx, \alpha gx) = d(gx, g\alpha x) = d(x, \alpha x) = \tau(x)
  \]
  as \(x\) is on the axis. Therefore \(g\) preserves \(\Axis(\alpha)\) so \(C(\alpha\) acts on \(\Axis(\alpha)\). We have \(\langle \alpha \rangle \subseteq C(\alpha)\). \(\alpha\) is injective then \(\deg p = 1 = |C(\alpha): \langle a \rangle|\) so the result follows.
\end{proof}

















\printindex
\end{document}

% Farb & Margalit: Mapping Class Groups